\documentclass[12pt,a4paper]{article}
\usepackage[margin=1in]{geometry}
\usepackage{graphicx}
\usepackage{hyperref}
\usepackage{booktabs}
\usepackage{listings}
\usepackage{float}
\title{\textbf{Smart Queue \& Crowd Prediction System}\\Project Proposal and Prototype}
\author{Student Name: \underline{\hspace{5cm}}\\Student ID: \underline{\hspace{5cm}}\\Department: Information and Communication Technology}
\date{\today}
\begin{document}
\maketitle
\tableofcontents
\newpage
\section{Introduction}
The Smart Queue \& Crowd Prediction System is a Java application designed to organize service queues,
issue token numbers, estimate waiting time, and show a simple crowd-level prediction.
\section{Problem Statement}
Manual queues can create uncertainty about a customer's position and expected waiting time.
This prototype demonstrates how a digital queue can make the process more organized.
\section{Objectives}
\begin{itemize}
\item Generate a token number for each customer during a running session.
\item Display customers currently waiting.
\item Allow an operator to call the next customer.
\item Estimate waiting time using a simple formula.
\item Display a rule-based crowd level: LOW, MEDIUM, or HIGH.
\end{itemize}
\section{Technology Stack}
\begin{itemize}
\item Programming language: Java
\item GUI: Java Swing
\item Storage: in-memory only in this prototype
\end{itemize}
\section{Classes and Responsibilities}
\begin{description}
\item[\texttt{Customer}] Stores customer name and requested service.
\item[\texttt{QueueToken}] Associates a token number with a customer.
\item[\texttt{QueueManager}] Adds customers, calls the next customer, and estimates waiting time.
\item[\texttt{CrowdPredictor}] Maps queue size to LOW, MEDIUM, or HIGH.
\item[\texttt{QueueDashboard}] Provides the graphical user interface.
\item[\texttt{SmartQueueSystem}] Starts the application.
\end{description}
\section{Waiting-Time Estimation}
The prototype estimates waiting time as
\[
W = \left\lceil \frac{N \times S}{C} \right\rceil,
\]
where \(N\) is the number of people ahead, \(S\) is average service time in minutes, and \(C\)
is the number of active counters. This is an estimate, not a guaranteed wait time.
\section{Crowd-Level Rules}
\begin{center}
\begin{tabular}{ll}
\toprule
Queue size & Crowd level \\
\midrule
0--5 & LOW \\
6--15 & MEDIUM \\
16 or more & HIGH \\
\bottomrule
\end{tabular}
\end{center}
\section{Software Process Model}
An iterative approach is suitable: implement token creation first, then queue display and calling,
then waiting-time estimation, and finally test and improve the GUI.
\section{Testing Plan}
\begin{itemize}
\item Add a customer with valid details.
\item Try to add a customer with missing details.
\item Verify that token numbers increase.
\item Call the next customer with a non-empty queue and with an empty queue.
\item Check waiting time with valid and invalid token numbers.
\item Verify crowd level as the queue size changes.
\end{itemize}
\section{Limitations and Future Work}
The current version has no database, authentication, SMS notifications, or machine-learning prediction.
Future versions can add MySQL persistence, admin login, QR-based tokens, notifications, and historical forecasting.
\section{How to Run}
From the project's \texttt{src} directory:
\begin{lstlisting}
javac *.java
java SmartQueueSystem
\end{lstlisting}
\section{Class Diagram}
Export the class diagram as \texttt{class-diagram.png} and place it beside this file if you wish to include it.
% \begin{figure}[H]
% \centering
% \includegraphics[width=0.95\linewidth]{class-diagram.png}
% \caption{Class diagram for the prototype}
% \end{figure}
\section{Conclusion}
This project demonstrates a digital queue workflow using Java Swing and six cooperating classes.
It is an academic prototype that can later be extended with persistent storage and additional features.
\end{document}
